\section{Limitations and threats to validity}
\label{sec:limitations}

\subsection{Synthetic data}

We use synthetic-but-realistic telemetry generated against Google's microservices-demo, with humanized Jira tickets calibrated against \textsc{TAWOS}'s 458{,}000 real Jira issues for length, comment count, and code-block ratio. We do not validate on production Jira from a real industrial team. Reviewers should treat our findings as characterizations of the \emph{retrieval-augmented architecture}, not predictions of production yields. We argue the architecture-level findings (depth scaling, the Hit@$K$ correction, channel ablation methodology, distractor robustness) generalize to other corpora, but the absolute Hit@1 = 0.250 number is corpus-specific.

\subsection{Single deployment}

All experiments use a single Google microservices-demo Kubernetes deployment with 14 fault families. Multi-deployment validation would strengthen the deployment-history scaling claim. We expect the qualitative result (monotone scaling with $n_{\text{prior}}$) to hold on different deployments because it follows from the structural property that retrieval is bounded by what's in memory; the quantitative magnitudes will vary with domain-specific factors (telemetry quality, ticket realism, retrieval-head architecture).

\subsection{Family-disjoint train/test split}

Train and test scenario families do not overlap. This is \emph{good} for out-of-distribution generalization claims --- the retrieval head cannot have memorized train-family specifics --- but it means the system's depth-scaling result is measured on families it has never seen at training time. We see this as a strength rather than a weakness: it shows the retrieval architecture generalizes structurally rather than relying on per-family supervision.

\subsection{V2 channel uniformity}

As reported in Section~\ref{sec:results} (RQ4), the V2 humanized corpus is dominated by engineer-voice log text and does not include explicit trace IDs or kubernetes nouns. The trace and k8s channel ablations are therefore uninformative on V2. Future humanizers should emit explicit multi-channel evidence (e.g., \texttt{traceId=abc123} lines, \texttt{Pod cartservice-7d4f}/\texttt{kubelet} mentions) so per-channel ablations carry signal. We disclose this explicitly rather than reporting the trace/k8s no-ops as if they were meaningful negatives.

\subsection{Novelty detection failure}

333 cold-start anomalies are not flagged correctly as novel. The current rule is a hard-coded score threshold; a learned classifier should do better. We did not implement and evaluate the learned classifier in this paper because doing so cleanly would require a separate set of training data (per-window novelty labels with sufficient class balance), and we deemed that scope too large for the current paper.

\subsection{Cross-encoder fine-tune size}

We fine-tune on 12{,}641 pairs, which is small by standard cross-encoder fine-tuning benchmarks (MS-MARCO has 39 million query-passage pairs). The modest +10--20\% relative Hit@5 improvement is consistent with this scale; richer training data (more positive pairs per window, or pairs from real production teams) should yield larger gains. We did not pursue a larger fine-tune because of corpus-size limits, not because of fundamental method constraints.

\subsection{Engineer-time-saved simulation assumptions}

The time-to-diagnose simulation (RQ6) assumes 30\,s per candidate reviewed and a 30-minute fallback. These numbers are our best estimates from informal industry conversation and are not measured against a real on-call team. The relative comparison across pipelines (HGB 30 min versus \sotaft{} 23.6 min) is more robust than the absolute numbers --- changing the assumptions linearly rescales all bars without changing the ordering or the per-pipeline ratios.

\subsection{Static memory snapshot}

We treat the memory corpus as a static snapshot at test time. In a real deployment, memory grows continuously as new incidents are resolved. We approximate this with the time-ordered visibility rule (a test window at time $t$ sees only tickets minted before $t$), but our depth-scaling axis collapses this temporal structure into a single ``$n_{\text{prior}}$'' count. A truly temporal evaluation would measure retrieval quality as a function of time-since-deployment, holding the test stream fixed. We argue our cross-sectional depth analysis is the right first step: it characterizes the steady-state behavior; future work can extend to the temporal-dynamics characterization.
